\documentclass[11pt,a4paper]{ctexart}

\usepackage[margin=2.2cm]{geometry}
\usepackage{amsmath,amssymb}
\usepackage{booktabs}
\usepackage{enumitem}
\usepackage{xcolor}
\usepackage{listings}
\usepackage{hyperref}
\usepackage{fancyhdr}
\usepackage{parskip}

\hypersetup{
  colorlinks=true,
  linkcolor=blue!50!black,
  urlcolor=blue!60!black
}

\pagestyle{fancy}
\fancyhf{}
\lhead{StatProg2}
\rhead{Practical 2: Functions}
\cfoot{\thepage}

\definecolor{codebg}{RGB}{247,247,247}
\definecolor{codegreen}{RGB}{0,110,60}
\definecolor{codepurple}{RGB}{120,40,120}
\definecolor{codegray}{RGB}{100,100,100}

\lstdefinelanguage{R}{
  morekeywords={
    function,if,else,for,while,repeat,in,next,break,TRUE,FALSE,NULL,
    NA,NaN,Inf,return,library
  },
  sensitive=true,
  morecomment=[l]{\#},
  morestring=[b]"
}

\lstset{
  language=R,
  backgroundcolor=\color{codebg},
  basicstyle=\ttfamily\small,
  keywordstyle=\color{blue!65!black}\bfseries,
  commentstyle=\color{codegreen},
  stringstyle=\color{codepurple},
  numberstyle=\tiny\color{codegray},
  numbers=left,
  numbersep=8pt,
  frame=single,
  rulecolor=\color{black!15},
  breaklines=true,
  showstringspaces=false,
  tabsize=2,
  columns=fullflexible
}

\setlist[itemize]{leftmargin=2em}
\setlist[enumerate]{leftmargin=2em}

\title{\textbf{StatProg2 -- Practical 2: Functions}\\
\large 题目、答案与知识点整理 / Aufgaben, Lösungen und Konzepte}
\author{}
\date{}

\begin{document}

\maketitle
\tableofcontents
\newpage

\section{Exercise 1: Naming Functions}

\subsection{Task 1: 给函数重新命名}

给定函数：

\begin{lstlisting}
f1 <- function(string, prefix) {
  str_sub(string, 1, str_length(prefix)) == prefix
}

f3 <- function(x, y) {
  rep(y, length.out = length(x))
}
\end{lstlisting}

\subsubsection*{答案}

第一个函数检查一个字符串是否以指定前缀开始，因此可以命名为：

\begin{lstlisting}
has_prefix <- function(string, prefix) {
  str_sub(string, 1, str_length(prefix)) == prefix
}
\end{lstlisting}

第二个函数把 \texttt{y} 重复到与 \texttt{x} 相同的长度，因此可以命名为：

\begin{lstlisting}
recycle_to_length <- function(x, y) {
  rep(y, length.out = length(x))
}
\end{lstlisting}

也可以使用更明确的参数名：

\begin{lstlisting}
recycle_to_length <- function(value, template) {
  rep(value, length.out = length(template))
}
\end{lstlisting}

\subsubsection*{知识点}

\begin{itemize}
  \item 好的函数名通常使用动词或动词短语，例如
  \texttt{has\_prefix()}、\texttt{calculate\_mean()}。
  \item 参数名应该说明参数的角色。与 \texttt{x}、\texttt{y} 相比，
  \texttt{string}、\texttt{prefix} 更容易理解。
  \item R 项目中通常使用 \texttt{snake\_case}。
\end{itemize}

\subsubsection*{Deutsch}

Die erste Funktion prüft, ob ein String mit einem bestimmten Präfix beginnt.
Ein sinnvoller Name ist daher \texttt{has\_prefix()}.

Die zweite Funktion recycelt einen Wert oder Vektor, bis die Länge eines
Referenzvektors erreicht ist. Ein geeigneter Name ist
\texttt{recycle\_to\_length()}.

\subsection{Task 2: \texttt{rnorm()} oder \texttt{norm\_r()}?}

R verwendet für Wahrscheinlichkeitsverteilungen ein einheitliches Schema:

\begin{center}
\begin{tabular}{lll}
\toprule
Präfix & Bedeutung & Normalverteilung \\
\midrule
\texttt{r} & Zufallswerte erzeugen & \texttt{rnorm()} \\
\texttt{d} & Dichte & \texttt{dnorm()} \\
\texttt{p} & kumulative Wahrscheinlichkeit & \texttt{pnorm()} \\
\texttt{q} & Quantil & \texttt{qnorm()} \\
\bottomrule
\end{tabular}
\end{center}

\subsubsection*{Warum könnte \texttt{norm\_r()} besser sein?}

\begin{itemize}
  \item Alle Funktionen der Normalverteilung würden mit \texttt{norm\_}
  beginnen.
  \item Durch Autovervollständigung wären
  \texttt{norm\_r()}, \texttt{norm\_d()}, \texttt{norm\_p()} und
  \texttt{norm\_q()} gemeinsam sichtbar.
\end{itemize}

\subsubsection*{Warum ist \texttt{rnorm()} trotzdem sinnvoller?}

\begin{itemize}
  \item Die Namen folgen einer einheitlichen R-Konvention.
  \item Das gleiche Schema gilt für andere Verteilungen:
\end{itemize}

\begin{lstlisting}
rbinom()
dbinom()
pbinom()
qbinom()

runif()
dunif()
punif()
qunif()
\end{lstlisting}

Wenn man die Präfixe \texttt{r}, \texttt{d}, \texttt{p} und \texttt{q}
kennt, kann man viele Funktionsnamen vorhersagen.

Noch explizitere Namen wären:

\begin{lstlisting}
normal_sample()
normal_density()
normal_cdf()
normal_quantile()
\end{lstlisting}

\section{Exercise 2: Funktionen schreiben und testen}

\subsection{Task 1: Penguin-Plot in eine Funktion verpacken}

\subsubsection*{Lösung}

\begin{lstlisting}
library(palmerpenguins)
library(tidyverse)

make_label <- function(variable_name) {
  str_replace_all(variable_name, "_", " ") %>%
    str_to_title()
}

scatter_penguins <- function(
  species_name,
  x_var = "bill_length_mm",
  y_var = "bill_depth_mm"
) {
  penguins %>%
    filter(species == .env$species_name) %>%
    ggplot(
      aes(
        x = .data[[x_var]],
        y = .data[[y_var]]
      )
    ) +
    geom_point() +
    labs(
      title = paste(
        "Scatter plot of",
        make_label(x_var),
        "vs",
        make_label(y_var),
        "for",
        species_name,
        "penguins"
      ),
      x = make_label(x_var),
      y = make_label(y_var)
    )
}
\end{lstlisting}

Aufruf:

\begin{lstlisting}
scatter_penguins("Adelie")

scatter_penguins(
  "Gentoo",
  "flipper_length_mm",
  "body_mass_g"
)
\end{lstlisting}

\subsubsection*{关键知识点}

\begin{itemize}
  \item \texttt{species} 是数据框中的列。
  \item \texttt{.env\$species\_name} 表示函数环境中的参数。
  \item \texttt{.data[[x\_var]]} 根据字符串形式的列名访问数据列。
  \item R 函数默认返回最后一个表达式，因此不一定需要
  \texttt{return(plot)}。
\end{itemize}

\subsubsection*{常见错误}

错误写法：

\begin{lstlisting}
aes(y = .data[[y = y_var]])
\end{lstlisting}

正确写法：

\begin{lstlisting}
aes(y = .data[[y_var]])
\end{lstlisting}

\texttt{.env} 只能在 \texttt{filter()}、\texttt{mutate()}、
\texttt{aes()} 等数据掩码环境中使用。在普通的 \texttt{paste()} 中应直接写：

\begin{lstlisting}
paste(..., species_name, ...)
\end{lstlisting}

另外，\texttt{"Torgersen"} 是 \texttt{island} 的取值，不是
\texttt{species}。企鹅物种是：

\begin{lstlisting}
Adelie
Chinstrap
Gentoo
\end{lstlisting}

\subsection{Task 2: Docstrings mit roxygen2}

\begin{lstlisting}
#' Convert a variable name into a readable label
#'
#' Replaces underscores with spaces and converts the string
#' to title case.
#'
#' @param variable_name A character string containing a
#'   variable name.
#' @return A readable character string.
#'
#' @examples
#' make_label("bill_length_mm")
make_label <- function(variable_name) {
  str_replace_all(variable_name, "_", " ") %>%
    str_to_title()
}


#' Create a scatter plot for one penguin species
#'
#' @param species_name A penguin species. Expected values are
#'   "Adelie", "Chinstrap", and "Gentoo".
#' @param x_var Name of the numeric x-axis variable.
#' @param y_var Name of the numeric y-axis variable.
#' @return A ggplot object.
#'
#' @examples
#' scatter_penguins("Adelie")
#' scatter_penguins(
#'   "Gentoo",
#'   "flipper_length_mm",
#'   "body_mass_g"
#' )
scatter_penguins <- function(
  species_name,
  x_var = "bill_length_mm",
  y_var = "bill_depth_mm"
) {
  penguins %>%
    filter(species == .env$species_name) %>%
    ggplot(
      aes(
        x = .data[[x_var]],
        y = .data[[y_var]]
      )
    ) +
    geom_point() +
    labs(
      title = paste(
        make_label(x_var),
        "vs.",
        make_label(y_var),
        "for",
        species_name,
        "penguins"
      ),
      x = make_label(x_var),
      y = make_label(y_var)
    )
}
\end{lstlisting}

\subsubsection*{Deutsch}

Roxygen2-Dokumentation beginnt mit \texttt{\#'}.
\texttt{@param} beschreibt Argumente, \texttt{@return} den Rückgabewert
und \texttt{@examples} typische Aufrufe.

\subsection{Task 3: Randfälle testen}

\begin{lstlisting}
scatter_penguins("Penguin")
scatter_penguins("Adelie", x_var = "fluffiness")
scatter_penguins("Adelie", x_var = "island")
\end{lstlisting}

\begin{center}
\begin{tabular}{p{3.5cm}p{5.8cm}p{5.2cm}}
\toprule
Test & Verhalten vor Validierung & Problem \\
\midrule
Ungültige Art & leerer Plot, meist kein Fehler &
stiller Fehler \\
Nicht vorhandene Spalte & Fehler aus \texttt{ggplot2} &
Fehlermeldung ist relativ technisch \\
Nicht numerische Spalte & Plot mit diskreter Achse &
technisch gültig, aber gegen die Funktionsidee \\
\bottomrule
\end{tabular}
\end{center}

\subsubsection*{中文说明}

\begin{itemize}
  \item 不存在的物种产生 0 行数据，因此通常只得到空图。
  \item 不存在的列会在 \texttt{ggplot2} 计算映射时出错。
  \item 分类变量可以被 \texttt{ggplot2} 当作离散轴，因此不一定报错。
\end{itemize}

\subsection{Task 4: Hilfsfunktionen verwenden}

\begin{lstlisting}
make_label <- function(variable_name) {
  str_replace_all(variable_name, "_", " ") %>%
    str_to_title()
}

filter_penguin_species <- function(data, species_name) {
  data %>%
    filter(species == .env$species_name)
}

scatter_plot <- function(data, x_var, y_var, title) {
  ggplot(
    data,
    aes(
      x = .data[[x_var]],
      y = .data[[y_var]]
    )
  ) +
    geom_point() +
    labs(
      title = title,
      x = make_label(x_var),
      y = make_label(y_var)
    )
}

scatter_penguins <- function(
  species_name,
  x_var = "bill_length_mm",
  y_var = "bill_depth_mm"
) {
  title <- paste(
    "Scatter plot of",
    make_label(x_var),
    "vs",
    make_label(y_var),
    "for",
    species_name,
    "penguins"
  )

  penguins %>%
    filter_penguin_species(species_name) %>%
    scatter_plot(
      x_var = x_var,
      y_var = y_var,
      title = title
    )
}
\end{lstlisting}

\subsubsection*{知识点}

每个辅助函数只负责一个任务：

\begin{center}
\begin{tabular}{ll}
\toprule
Funktion & Aufgabe \\
\midrule
\texttt{make\_label()} & lesbare Achsenbeschriftung \\
\texttt{filter\_penguin\_species()} & Daten filtern \\
\texttt{scatter\_plot()} & Scatterplot erstellen \\
\texttt{scatter\_penguins()} & einzelne Schritte koordinieren \\
\bottomrule
\end{tabular}
\end{center}

\section{Exercise 3: Input Validation}

\subsection{Task 1: Eingaben prüfen}

\begin{lstlisting}
scatter_penguins <- function(
  species_name,
  x_var = "bill_length_mm",
  y_var = "bill_depth_mm"
) {
  valid_species <- c("Adelie", "Chinstrap", "Gentoo")

  if (!species_name %in% valid_species) {
    stop(
      "`species_name` must be one of: ",
      paste(valid_species, collapse = ", "),
      ". You supplied: \"",
      species_name,
      "\"."
    )
  }

  stopifnot(
    "`x_var` must be a column of `penguins`" =
      x_var %in% names(penguins),

    "`y_var` must be a column of `penguins`" =
      y_var %in% names(penguins),

    "`x_var` must be numeric" =
      is.numeric(penguins[[x_var]]),

    "`y_var` must be numeric" =
      is.numeric(penguins[[y_var]])
  )

  penguins %>%
    filter(species == .env$species_name) %>%
    ggplot(
      aes(
        x = .data[[x_var]],
        y = .data[[y_var]]
      )
    ) +
    geom_point() +
    labs(
      title = paste(
        make_label(x_var),
        "vs.",
        make_label(y_var),
        "for",
        species_name,
        "penguins"
      ),
      x = make_label(x_var),
      y = make_label(y_var)
    )
}
\end{lstlisting}

\subsubsection*{检查顺序}

必须先检查列是否存在，再检查列是否为数值：

\[
\text{Spalte vorhanden?}
\longrightarrow
\text{Datentyp numerisch?}
\]

\subsubsection*{Testfälle}

\begin{lstlisting}
scatter_penguins("Emperor")
# Fehler: ungueltige Art

scatter_penguins("Adelie", "fluffiness", "body_mass_g")
# Fehler: x_var ist keine Spalte

scatter_penguins("Adelie", "island", "body_mass_g")
# Fehler: x_var muss numerisch sein
\end{lstlisting}

\subsection{Task 2: Alte und neue Fehlermeldungen vergleichen}

Die validierte Version ist leichter zu debuggen:

\begin{itemize}
  \item Sie scheitert früh (\emph{fail fast}).
  \item Sie nennt direkt das fehlerhafte Argument.
  \item Sie verhindert stille Fehler wie leere Plots.
  \item Sie erklärt, welche Eingaben zulässig sind.
\end{itemize}

\subsubsection*{Musterantwort}

\begin{quote}
The validated version is easier to debug because it fails early, uses the
names of the function arguments, avoids silent failures, and tells the user
how to correct the input.
\end{quote}

\section{Exercise 4: Iteration und Rekursion}

\subsection{Task 1: Fibonacci iterativ}

Die Fibonacci-Folge ist definiert durch:

\[
F(0)=0,\qquad F(1)=1,
\]

und für \(n>1\):

\[
F(n)=F(n-1)+F(n-2).
\]

\begin{lstlisting}
fib_iterative <- function(n) {
  if (n == 0) {
    return(0)
  } else if (n == 1) {
    return(1)
  }

  a <- 0
  b <- 1

  for (i in 2:n) {
    temp <- a + b
    a <- b
    b <- temp
  }

  return(b)
}
\end{lstlisting}

Beispiele:

\begin{lstlisting}
fib_iterative(0)
# 0

fib_iterative(5)
# 5

fib_iterative(10)
# 55
\end{lstlisting}

\subsubsection*{中文解释}

初始时：

\[
a=F(0)=0,\qquad b=F(1)=1.
\]

每一轮先计算：

\[
\texttt{temp}=a+b,
\]

然后把两个保存值向后移动：

\[
a\leftarrow b,\qquad b\leftarrow \texttt{temp}.
\]

使用 \texttt{temp} 是为了避免旧的 \texttt{a} 在加法前被覆盖。

\subsection{Task 2: Fibonacci rekursiv}

\begin{lstlisting}
fib_recursive <- function(n) {
  if (n == 0) {
    return(0)
  } else if (n == 1) {
    return(1)
  }

  return(
    fib_recursive(n - 1) +
      fib_recursive(n - 2)
  )
}
\end{lstlisting}

递归函数包含：

\begin{enumerate}
  \item 基本情况：\texttt{n == 0} 和 \texttt{n == 1}；
  \item 递归情况：调用 \texttt{fib\_recursive(n - 1)} 和
  \texttt{fib\_recursive(n - 2)}。
\end{enumerate}

Für \(F(5)\):

\[
F(5)=F(4)+F(3)=3+2=5.
\]

Ohne Basisfälle würde sich die Funktion unbegrenzt weiter aufrufen.

\subsection{Task 3: Laufzeit vergleichen}

\begin{lstlisting}
system.time(fib_iterative(30))
system.time(fib_recursive(30))
\end{lstlisting}

Beide Funktionen liefern:

\[
F(30)=832040.
\]

Die iterative Variante ist aber wesentlich schneller.

\begin{center}
\begin{tabular}{lll}
\toprule
Eigenschaft & Iterativ & Naiv rekursiv \\
\midrule
Berechnung & jedes Ergebnis einmal & viele Wiederholungen \\
Laufzeit & \(O(n)\) & ungefähr \(O(\varphi^n)\) \\
Zusatzspeicher & \(O(1)\) & Aufrufstapel \(O(n)\) \\
Lesbarkeit & etwas länger & nahe an der Formel \\
\bottomrule
\end{tabular}
\end{center}

Die Rekursion ist nicht grundsätzlich langsam. Das Problem ist hier, dass
dieselben Teilprobleme mehrfach berechnet werden. Memoization könnte diese
Wiederholungen vermeiden.

\section{Exercise 5: Refactoring}

\subsection{Task 1.1: Wiederholungen in \texttt{summarise\_species()}}

Die ursprüngliche Funktion wiederholt:

\begin{itemize}
  \item die Filterbedingung für jede Art,
  \item die Berechnung von Mittelwerten,
  \item \texttt{mean(..., na.rm = TRUE)},
  \item die Ausgabe mit \texttt{cat()},
  \item lange \texttt{if}/\texttt{else if}-Zweige.
\end{itemize}

Trotz ihrer Länge enthält sie nur drei Ideen:

\begin{enumerate}
  \item auf eine Art filtern,
  \item mittlere Schnabellänge und Körpermasse berechnen,
  \item Ergebnisse ausgeben.
\end{enumerate}

\subsection{Task 1.2: \texttt{summarise\_species()} refaktorieren}

\begin{lstlisting}
summarise_species <- function(
  species_name,
  data = penguins
) {
  summary_row <- data %>%
    filter(species == .env$species_name) %>%
    summarise(
      mean_bill = mean(
        bill_length_mm,
        na.rm = TRUE
      ),
      mean_mass = mean(
        body_mass_g,
        na.rm = TRUE
      )
    )

  cat(
    species_name,
    "-- mean bill:",
    summary_row$mean_bill,
    "mean mass:",
    summary_row$mean_mass,
    "\n"
  )
}
\end{lstlisting}

Aufrufe:

\begin{lstlisting}
summarise_species("Adelie")
summarise_species("Gentoo")
summarise_species("Chinstrap")
\end{lstlisting}

\subsubsection*{Warum funktioniert eine vierte Art?}

Die Funktion enthält keine fest codierten Berechnungen für einzelne Arten.
Sie verwendet direkt \texttt{species\_name}. Wenn der Datensatz später eine
weitere Art enthält, kann sie ohne neuen \texttt{else if}-Zweig verarbeitet
werden.

Robustere Variante:

\begin{lstlisting}
summarise_species <- function(
  species_name,
  data = penguins
) {
  valid_species <- unique(as.character(data$species))

  if (!species_name %in% valid_species) {
    stop(
      "`species_name` must be one of: ",
      paste(valid_species, collapse = ", ")
    )
  }

  summary_row <- data %>%
    filter(species == .env$species_name) %>%
    summarise(
      mean_bill = mean(
        bill_length_mm,
        na.rm = TRUE
      ),
      mean_mass = mean(
        body_mass_g,
        na.rm = TRUE
      )
    )

  cat(
    species_name,
    "-- mean bill:",
    summary_row$mean_bill,
    "mean mass:",
    summary_row$mean_mass,
    "\n"
  )

  invisible(summary_row)
}
\end{lstlisting}

\subsection{Task 2: \texttt{penguin\_report()} in Hilfsfunktionen zerlegen}

\begin{lstlisting}
load_penguins <- function() {
  palmerpenguins::penguins
}

drop_missing <- function(data, cols) {
  data[complete.cases(data[, cols]), ]
}

print_summary <- function(data) {
  cat("n =", nrow(data), "\n")
  cat(
    "mean bill length:",
    mean(data$bill_length_mm),
    "\n"
  )
  cat(
    "mean body mass:",
    mean(data$body_mass_g),
    "\n"
  )
}

scatter_bill_vs_mass <- function(data, title) {
  ggplot(
    data,
    aes(
      x = bill_length_mm,
      y = body_mass_g
    )
  ) +
    geom_point() +
    labs(title = title)
}

penguin_report <- function(species_name) {
  data <- load_penguins() %>%
    drop_missing(
      c(
        "bill_length_mm",
        "body_mass_g"
      )
    ) %>%
    filter(species == .env$species_name)

  print_summary(data)

  scatter_bill_vs_mass(
    data,
    title = paste(
      "Bill length vs. body mass --",
      species_name
    )
  )
}
\end{lstlisting}

Aufruf:

\begin{lstlisting}
penguin_report("Gentoo")
\end{lstlisting}

\subsubsection*{Vorteile der Refaktorierung}

\begin{itemize}
  \item Jede Funktion übernimmt nur eine Aufgabe.
  \item Einzelne Schritte können separat getestet werden.
  \item Reinigungs- und Plotfunktionen sind wiederverwendbar.
  \item Die Hauptfunktion liest sich wie ein Rezept:
\end{itemize}

\[
\text{laden}
\rightarrow
\text{bereinigen}
\rightarrow
\text{filtern}
\rightarrow
\text{zusammenfassen}
\rightarrow
\text{visualisieren}.
\]

\section{Zusammenfassung der wichtigsten Konzepte}

\begin{enumerate}
  \item Funktionen sollten klare Namen und verständliche Parameter besitzen.
  \item Wiederholte Codeblöcke sollten durch Funktionen ersetzt werden.
  \item Veränderliche Teile gehören in Funktionsargumente.
  \item \texttt{.data[[name]]} greift auf eine Spalte zu, deren Name als
  String gespeichert ist.
  \item \texttt{.env\$name} greift innerhalb eines Data-Masking-Kontexts
  auf ein Funktionsargument zu.
  \item Eingaben sollten früh geprüft werden.
  \item Gute Fehlermeldungen nennen das Argument, das Problem und möglichst
  die erlaubten Werte.
  \item Rekursion benötigt einen Basisfall.
  \item Naive Fibonacci-Rekursion ist wegen mehrfacher Berechnung langsam.
  \item Große Funktionen sollten in kleine Hilfsfunktionen mit klarer
  Verantwortung zerlegt werden.
\end{enumerate}

\end{document}
